\begin{description}
  \item [Dupicate finder]: A class that finds duplicate images

    This class can be configured the pHash or dHash algorithm for fuzzy hashing. It exposes 
    a method called \texttt{find\_duplicates} expects a list of filepaths as an input. The 
    given files are then loaded in batches, hashed and stored in a Dictionary if the hamming
    distance is higher than a given threshold. The Dictionary is used to avoid unnecessary 
    comparisons (images are in the same "Bucket if the first 32 Bits match"). The method returns 
    a list of paths to duplicte images for further processing.

  \item [Blur filter]: A class that measures the blurriness of a given image

    This class is part of the group's preprocessing pipeline. Given an image the class computes
    the laplacian variance of the input and removes it if the variance is lower than a given threshold. 
    This implementation was later adapted to use the \textit{fast fourier transform} (FFT) to classify the 
    image.

  \item [Different Model implementations]: Different custom implementations that were trained

    During the development phase of the project I wanted to evaluate as many different model architectures 
    as possible. Due to the fixed input size of ($64 \times 64$) all models were implemented from scratch 
    and modified to downsample less aggressively. The developed models were all trained using a custom training 
    script.

  \item [Model Evaluation script]: A script that given a Model implementation measures its accuracy

    This script was later used by all members of the group to measure and compare their model's performance on 
    unseen datasets.
  
  \item [Image demo]: A script that processes all images in a directory

    This script expects the path to a directory filled with images, loads all images in batches,
    loads our model and evaluates all images. The resulting probabilites for each emotion are then
    written into a csv file with the related filepath


\end{description}
